\section*{ID10259: Entropic Gravity Emergence}
\paragraph{Metadata.}
PiID: 4077365949563 | RTEID: RTE:P39837.9753:T2.8089 | UUID: 259935b9-475b-5798-a048-e6d9074efd24

\paragraph{Canonical Statement.}
\textbackslash{}begin\{proof\} Apply Verlinde's entropic force principle with holographic screen area \$A = 4\textbackslash{}pi r\_h\textasciicircum{}2\$: \textbackslash{}begin\{equation\} F\_\{\textbackslash{}text\{entropic\}\} = T \textbackslash{}nabla S = \textbackslash{}frac\{k\_B T\}\{2\textbackslash{}pi r\_h\} \textbackslash{}cdot \textbackslash{}frac\{A c\textasciicircum{}3\}\{4G\textbackslash{}hbar\} = \textbackslash{}frac\{GMm\}\{r\textasciicircum{}2\} \textbackslash{}left(1 + \textbackslash{}mathcal\{O\}\textbackslash{}left(\textbackslash{}frac\{\textbackslash{}ell\_P\}\{r\}\textbackslash{}right)\textbackslash{}right) \textbackslash{}end\{equation\} Temporal integration via Duhamel's principle yields the accumulated tensor. \textbackslash{}end\{proof\}

\paragraph{Canonical Equation.}
\[
A = 4\pi r_h^2
\]

\paragraph{Dictionary.}
Temporal integration via Duhamel's principle yields the accumulated tensor [ID10259]

\paragraph{Encyclopedia.}
Temporal integration via Duhamel's principle yields the accumulated tensor. \{proof\} It is an algebraic definition expressing the quantity directly in terms of the variables on its right-hand side.
